\chapter{Market Position Analysis / Marktpositions-Analyse / 市场地位分析}

\section{Market Share and Indicators / Marktanteil und Indikatoren / 市场份额与指标}

\begin{otherlanguage}{ngerman}
Edekas Marktposition:
\begin{itemize}
    \item \textbf{Marktanteil}: 26\% im deutschen Lebensmitteleinzelhandel
    \item \textbf{Filialanzahl}: Über 11.000
    \item \textbf{Umsatz}: 60+ Milliarden Euro
    \item \textbf{Wachstum}: [Zu aktualisieren]
    \item \textbf{Rentabilität}: [Zu aktualisieren]
\end{itemize}
\end{otherlanguage}

Edeka's market position:
\begin{itemize}
    \item \textbf{Market share}: 26\% in German food retail
    \item \textbf{Store count}: Over 11,000
    \item \textbf{Revenue}: 60+ billion euros
    \item \textbf{Growth}: [To be updated]
    \item \textbf{Profitability}: [To be updated]
\end{itemize}

艾德卡的市场地位：
\begin{itemize}
    \item \textbf{市场份额}：德国食品零售市场的26\%
    \item \textbf{门店数量}：超过11,000家
    \item \textbf{收入}：600多亿欧元
    \item \textbf{增长}：[待更新]
    \item \textbf{盈利能力}：[待更新]
\end{itemize}

\section{Key Performance Indicators / Wichtige Leistungskennzahlen / 关键绩效指标}

\begin{otherlanguage}{ngerman}
KPIs für Edeka:
\begin{itemize}
    \item Umsatz pro Quadratmeter
    \item Kundenfrequenz
    \item Durchschnittlicher Warenkorbwert
    \item Kundenzufriedenheit
    \item Mitarbeiterproduktivität
    \item Markenbekanntheit
\end{itemize}
\end{otherlanguage}

% Revenue Trend Chart
\begin{figure}[h]
\centering
\begin{tikzpicture}[scale=1.0]
    % Axes
    \draw[->, thick] (0,0) -- (10,0) node[right] {Year / Jahr};
    \draw[->, thick] (0,0) -- (0,7) node[above, rotate=90] {Revenue (Billion €) / Umsatz (Mrd. €)};
    
    % Grid
    \foreach \x in {1,2,3,4,5,6,7,8,9}
        \draw[gray!20] (\x,0) -- (\x,7);
    \foreach \y in {1,2,3,4,5,6}
        \draw[gray!20] (0,\y) -- (10,\y);
    
    % Data points (2019-2023)
    \coordinate (2019) at (1,5.85);
    \coordinate (2020) at (3,5.98);
    \coordinate (2021) at (5,6.12);
    \coordinate (2022) at (7,6.25);
    \coordinate (2023) at (9,6.38);
    
    % Draw line
    \draw[blue!70, thick, smooth] (2019) -- (2020) -- (2021) -- (2022) -- (2023);
    
    % Draw points
    \foreach \point in {(2019), (2020), (2021), (2022), (2023)} {
        \fill[blue!70] \point circle (3pt);
    }
    
    % Labels on points
    \node[above] at (2019) {58.5};
    \node[above] at (2020) {59.8};
    \node[above] at (2021) {61.2};
    \node[above] at (2022) {62.5};
    \node[above] at (2023) {63.8};
    
    % X-axis labels
    \node[below] at (1,0) {2019};
    \node[below] at (3,0) {2020};
    \node[below] at (5,0) {2021};
    \node[below] at (7,0) {2022};
    \node[below] at (9,0) {2023};
    
    % Y-axis labels
    \node[left] at (0,1) {55};
    \node[left] at (0,3) {58};
    \node[left] at (0,5) {61};
    \node[left] at (0,7) {64};
    
    % Title
    \node[above, font=\bfseries] at (5,7.5) {Revenue Trend / Umsatztrend / 收入趋势};
    
    % Fill area under curve
    \fill[blue!20] (2019) -- (2020) -- (2021) -- (2022) -- (2023) -- (9,0) -- (1,0) -- cycle;
\end{tikzpicture}
\caption{Revenue Trend (2019-2023) / Umsatztrend / 收入趋势}
\end{figure}

KPIs for Edeka:
\begin{itemize}
    \item Revenue per square meter
    \item Customer frequency
    \item Average basket value
    \item Customer satisfaction
    \item Employee productivity
    \item Brand awareness
\end{itemize}

艾德卡的KPI：
\begin{itemize}
    \item 每平方米收入
    \item 客户频率
    \item 平均购物篮价值
    \item 客户满意度
    \item 员工生产力
    \item 品牌知名度
\end{itemize}

\section{Market Trends Impact / Auswirkungen von Markttrends / 市场趋势影响}

\begin{otherlanguage}{ngerman}
\begin{itemize}
    \item E-Commerce-Wachstum
    \item Nachhaltigkeitstrends
    \item Preisbewusstsein der Verbraucher
    \item Convenience-Orientierung
    \item Gesundheitsbewusstsein
\end{itemize}
\end{otherlanguage}

\begin{itemize}
    \item E-commerce growth
    \item Sustainability trends
    \item Consumer price consciousness
    \item Convenience orientation
    \item Health consciousness
\end{itemize}

\begin{itemize}
    \item 电子商务增长
    \item 可持续性趋势
    \item 消费者价格意识
    \item 便利性导向
    \item 健康意识
\end{itemize}
